%% main.tex
%% Copyright 2026 Marcelo Fort Muñoz
%
% This work is provided without any warranty, express or implied.
% The author makes no representations regarding its accuracy,
% reliability, safety, or suitability for any particular purpose.
%
% This work is not intended, designed, tested, or certified for use
% in safety-critical applications or in any situation where its
% failure could result in property damage, personal injury, or loss
% of life. It must not be relied upon in such circumstances.
%
% Anyone choosing to use this work in such circumstances assumes
% responsibility for independently assessing its suitability and
% for any necessary verification and safeguards, to the extent
% permitted by applicable law.
%
% This work may be distributed and/or modified under the
% conditions of the LaTeX Project Public License, either version 1.3
% of this license or (at your option) any later version.
% The latest version of this license is in
%   https://www.latex-project.org/lppl.txt
% and version 1.3c or later is part of all distributions of LaTeX
% version 2008 or later.
%
% This work has the LPPL maintenance status `maintained'.
%
% The Current Maintainer of this work is Marcelo Fort Muñoz.
% Reachable at: marcelo.fm@member.fsf.org
%
% Please report any issues to the maintainer.
%
% The files covered by this license are listed in
% "manifest.txt" in the root of this distribution.
% Files not listed therein are not included in this work.
% Their respective copyright and licensing terms remain applicable.
%
% SPDX-License-Identifier: LPPL-1.3c
%%


\documentclass[11pt,a4paper]{article}
\usepackage[spanish]{babel}
\usepackage{csquotes}
\usepackage{fontspec}
\usepackage{geometry}
\geometry{left=2cm, right=2cm, top=2cm, bottom=2cm}
\usepackage{graphicx}
\usepackage{eso-pic}
\usepackage{microtype}
\usepackage{longtable}
\usepackage{multirow}




\usepackage[svgnames]{xcolor}
\usepackage[colorlinks=true, linkcolor=Indigo, urlcolor=Indigo]{hyperref}


\usepackage{fancyhdr}
\pagestyle{fancy} % Enables the custom headers and footers

\usepackage{caption}
\usepackage[backend=biber, style=apa]{biblatex}
%\DeclareLanguageMapping{spanish}{spanish-apa}


\usepackage{listings}

% Definición de estilo para Modelo Relacional
\lstdefinelanguage{ModeloRelacional}{
    keywords=[1]{CE, REFERENCIA},
    keywordstyle=[1]\color{blue}\bfseries,
    keywords=[2]{CP},
    keywords=[3]{PROVEEDOR, EXISTENCIA, LOCALIZACION, PRODUCTO, MATERIALPROMOCIONAL, MERCHANDISING, MOVIMIENTO, LINEA, VENTA, USUARIO, EVENTO},
    keywordstyle=[2]\color{red}\bfseries,
    sensitive=true,
    comment=[l]{\#},
    commentstyle=\color{gray}\itshape,
    basicstyle=\ttfamily\small,
    breaklines=true,
    breakatwhitespace=true,
    showstringspaces=false,
    tabsize=2,
    frame=single,
    framesep=5pt,
    rulecolor=\color{black!30},
    backgroundcolor=\color{gray!5}
%    morekeywords={\textcolor{red}{\textbf{[A-Z]+}}}
}

\lstset{
    language=ModeloRelacional,
    inputencoding=utf8,
    extendedchars=true,
    literate={á}{{\'a}}1 {é}{{\'e}}1 {í}{{\'i}}1 {ó}{{\'o}}1 {ú}{{\'u}}1
        {Á}{{\'A}}1 {É}{{\'E}}1 {Í}{{\'I}}1 {Ó}{{\'O}}1 {Ú}{{\'U}}1
        {ñ}{{\~n}}1 {Ñ}{{\~N}}1 {→}{{$\rightarrow$}}1
}

\fancyfoot[C]{Departamento de ingeniería informática - Memoria prácticas en empresa}              % Center
\fancyfoot[R]{\thepage}

\renewcommand{\headrulewidth}{0pt}

\addbibresource{biblio/practicas.bib}


\renewcommand{\baselinestretch}{1.5}

\setmainfont{Liberation Sans}
\renewcommand{\familydefault}{\sfdefault}

\begin{document}

    \begin{titlepage}

% --- IMAXE DE FONDO / DEREITA ---
        \AddToShipoutPictureBG*{
            \AtPageLowerLeft{
                \put(140, 170){%
                % 15.5cm obriga á imaxe a saír por fóra do folio
                    \includegraphics[width=14cm, keepaspectratio, trim=0 0 1 0, clip]{recursos/nebrija_2}
                }
            }
        }

% --- IMAXE DE ABAIXO DEREITA ---
        \AddToShipoutPictureBG*{
            \AtPageLowerLeft{
                \put(335, 35){
                    \includegraphics[width=9cm,keepaspectratio]{recursos/nebrija_pie}
                }
            }
        }

% --- BLOQUE SUPERIOR ESQUERDO ---
        \noindent
        \begin{minipage}[t]{1\textwidth}
            \vspace{0.5cm}
            {\Large \textcolor{gray}{\textbf{Escuela Politécnica Superior}}} \\[0.8cm]
            {\small \textbf{Nombre del alumno:} Marcelo Fort Muñoz \{mfortm\}@alumnos.nebrija.es\\}
            {\small \textbf{N.º de precolegiación CPETIG:} 0262\\}
            {\small \textbf{Nombres de los tutores profesional y académico:} Tobias Diekershoff y Adolfo Abalo Cascallar\\}
            {\small \textbf{Tipología de memoria:} Desarrollo}
        \end{minipage}

        \vspace{7.0cm}

% --- BLOQUE CENTRAL ESQUERDO
        \bgroup
        \setlength{\fboxsep}{0pt}
        \noindent\hspace*{-3.2cm}
        \colorbox{red}{%
        % A4=21cm de ancho. 21/2 (á mitade) = 10.5cm.
            \begin{minipage}[c][7.86cm][c]{11.11cm}
                \centering
                \begin{minipage}{8.5cm}
                    \centering
                    {\Huge \textcolor{white}{\textbf{Título del trabajo que se presenta}}}
                \end{minipage}
            \end{minipage}%
        }
        \egroup

        \vspace{2.5cm}

% --- BLOQUE INFERIOR ESQUERDO (Asignatura e Grao) ---
        \noindent\hspace{-0.75cm}
        \begin{minipage}[t]{\textwidth}
            \flushleft
            {\Large \textcolor{red}{\textbf{Evaluación de Capacidades \\ en la Empresa I}}} \\[0.3cm]
            {\textbf{Grado en Ingeniería Informática}}
        \end{minipage}

    \end{titlepage}

    \section*{Agradecimientos}\label{sec:agradecimientos}

    \clearpage

    D. Marcelo Fort Muñoz autoriza a que el presente trabajo se guarde y custodie en los repositorios de la Universidad Nebrija y además SI/NO autoriza a su disposición en abierto.
    \pagebreak

    \tableofcontents
    \clearpage

    \listoffigures
    \clearpage

    \listoftables
    \clearpage



    \section{Resumen}\label{sec:resumen}
    Este documento recoge las prácticas realizadas en la \textit{Free Software Foundation Europe} (FSFE),
    organización dedicada al fomento y la protección del software libre y abierto.
    El periodo permitió participar en el mantenimiento de la infraestructura tecnológica de la fundación,
    así como en el desarrollo y actualización de herramientas utilizadas por su equipo y la comunidad.

    Entre las principales tareas se incluyeron la mejora de los sistemas de protección de servidores,
    la actualización de servicios y dependencias, la migración de una API XML-RPC a una API REST,
    el despliegue de un servicio de gestión de podcasts y la adaptación de formularios y funcionalidades
    de la página web. También se desarrolló una herramienta de análisis estadístico de eventos y noticias
    y se elaboró documentación técnica para la configuración y el uso de servicios internos.

    Las prácticas fueron una oportunidad para aplicar y ampliar conocimientos de Linux, administración
    de sistemas, desarrollo web, automatización y seguridad informática. Asimismo, contribuyeron a reforzar
    competencias profesionales como el trabajo con metodologías ágiles, la coordinación con equipos
    multidisciplinares y la gestión de incidencias en entornos reales.

    \section{Prácticas}\label{sec:practicas}
    En esta sección se presentan los detalles de tanto la experiencia profesional en las prácticas,
los antecedentes personales y académicos, así como la descripción de la empresa y su funcionamiento.
\subsection{Antecedentes}\label{subsec:antecedentes}
%todo
\subsubsection{Entrevistas realizadas}\label{subsubsec:entrevistas}
%todo
\subsubsection{Información previa}\label{subsubsec:informacion-previa}
%todo
\subsubsection{Motivos de la elección de prácticas}\label{subsubsec:motivos-eleccion}
%todo
\subsubsection{Datos de la empresa}\label{subsubsec:datos-empresa}
%todo
\subsection{Objetivos de las prácticas}\label{subsec:objetivos-practicas}
%todo
\subsection{Onboarding}\label{subsec:onboarding}
%todo
\subsection{Funciones y tareas en las prácticas}\label{subsec:funciones-tareas}
%todo
\subsection{Herramientas utilizadas en las las prácticas}\label{subsec:herramientas}
%todo
\subsection{Metodologías utilizadas}\label{subsec:metodologias}
%todo
\subsection{Valoración personal de la elección de las prácticas}\label{subsec:valoracion-personal}
%todo
%FIN SECCIÓN


    \section{Desarrollo}\label{sec:desarrollo}
    %% desarrollo.tex
%% Copyright 2026 Marcelo Fort Muñoz
%
% This work may be distributed and/or modified under the
% conditions of the LaTeX Project Public License, either version 1.3
% of this license or (at your option) any later version.
% The latest version of this license is in
%   https://www.latex-project.org/lppl.txt
% and version 1.3c or later is part of all distributions of LaTeX
% version 2008 or later.
%
% This work has the LPPL maintenance status `maintained'.
%
% The Current Maintainer of this work is Marcelo Fort Muñoz.
% This work consists of the files main.tex, secciones/practicas.tex, secciones/resumen.tex,
% biblio/ejemplo.bib, and any other files in the same directory (src/ and subfolders with the exception of png files).
% SPDX-License-Identifier: LPPL-1.3c

\subsection{Introducción}\label{subsec:introduccion}
\subsection{Antecedentes}\label{subsec:antecedentes}
\subsection{Analisis de mercado}\label{subsec:analisis-mercado}
\subsection{Necesidad detectada}\label{subsec:necesidad-detectada}
\subsection{Objetivos}\label{subsec:objetivos}


    \section{Marco teórico}\label{sec:marco-teorico}

    \section{Metodología}\label{sec:metodologia}

    \section{Proyecto}\label{sec:proyecto}
    \subsection{Resumen de contribuciones y productos desarrollados}\label{subsec:resumen-contribuciones}
\subsection{Planificación temporal}\label{subsec:planificacion-temporal}
\subsection{Recursos empleados}\label{subsec:recursos-empleados}
\subsection{Valoración económica}\label{subsec:valoracion-economica}


    \section{Resultados y discusión}\label{sec:resultados-discusion}
    Semántica no reflejada:
\begin{enumerate}
    \item \textbf{Cuando se genera un movimiento de venta con evento, el evento se guardará exclusivamente en la tabla de ventas, y no en la tabla de eventos.}
    \item Cantidad de movimiento no puede ser negativa ni superior a la cantidad de stock modificado.
\end{enumerate}

\begin{lstlisting}[caption={Modelo Relacional - Papertory}, label={fig:modelo-relacional-papertory}]

    USUARIO (id_usuario, nombre, email, password, Rol)
        CP: id_usuario

    PROVEEDOR (id_proveedor, nombre, email, vatId, telefono, direccion, ciudad, pais)
        CP: id_proveedor

    LOCALIZACION (id_localizacion, nombre)
        CP: id_localizacion

    PRODUCTO (id_producto, nombre, descripcion, codigo, activo, tipo)
        CP: id_producto

    MATERIALPROMOCIONAL (id_producto, tipo, ancho, alto, descripcion)
        CP: id_producto
        CE: id_producto REFERENCIA PRODUCTO(id_producto)

    MERCHANDISING (id_producto, talla, precio, iva, descripcion)
        CP: id_producto
        CE: id_producto REFERENCIA PRODUCTO(id_producto)

    EXISTENCIA (id_existencia, id_producto, id_localizacion, id_proveedor, fecha_entrada, fecha_apertura, cantidad_inicial, cantidad_actual)
        CP: id_existencia
        CE: id_producto REFERENCIA PRODUCTO(id_producto)
        CE: id_localizacion REFERENCIA LOCALIZACION(id_localizacion)
        CE: id_proveedor REFERENCIA PROVEEDOR(id_proveedor)

    EVENTO (id_evento, nombre, descripcion)
        CP: id_evento

    MOVIMIENTO (id_movimiento, id_existencia, id_evento, id_usuario, fecha, cantidad, tipo)
        CP: id_movimiento
        CE: id_existencia REFERENCIA EXISTENCIA(id_existencia)
        CE: id_evento REFERENCIA EVENTO(id_evento)
        CE: id_usuario REFERENCIA USUARIO(id_usuario)

    VENTA (id_venta, id_evento, id_usuario, fecha, metodo)
        CP: id_venta
        CE: id_evento REFERENCIA EVENTO(id_evento)
        CE: id_usuario REFERENCIA USUARIO(id_usuario)

    LINEA (id_linea, id_venta, id_movimiento, iva, precio_unitario, cantidad)
        CP: id_linea, id_venta
        CE: id_venta REFERENCIA VENTA(id_venta)
        CE: id_movimiento REFERENCIA MOVIMIENTO(id_movimiento)



    
\end{lstlisting}


\begin{longtable}{|l|l|p{3.8cm}|p{2.2cm}|c|c|c|c|}
    \hline
    \multicolumn{8}{| c |}{Begin of Table} \\
    \hline
    Entidad o relación & Atributos & Descripción & Tipo de datos & N & M & D & P \\
    \hline
    \multirow{6}{*}{Persona}
    & Nombre & Nombre y apellidos de cada usuario & Carácter variable 100 & NO & NO & NO & NO \\
    \cline{2-8}
    & Id\_usuario \{PK\} & Identificador único de cada usuario & Carácter variable 30 & NO & NO & NO & NO \\
    \cline{2-8}
    & Correo & Dirección de email único de cada usuario & Carácter variable 100 & NO & NO & NO & NO \\
    \cline{2-8}
    & Contraseña & Contraseña de cada usuario para acceder a la aplicación & Carácter variable 20 & NO & NO & NO & NO \\
    \cline{2-8}
    & Rol & Rol de cada usuario (administrador, empleado, cliente) & Carácter variable 20 & NO & NO & NO & NO \\
    \cline{2-8}
    & tipo\_usuario & Tipo de usuario: administrador, usuario & Carácter variable 20 & NO & NO & NO & 'Normal' \\
    \hline
    \hline
    \caption{This is the caption for the test table.}
\end{longtable}\label{tab:diccionario-datos}


    \section{Conclusiones}\label{sec:conclusiones}

    \section{Líneas futuras}\label{sec:lineas-futuras}

    \section{Uso de IA}\label{sec:uso-ia}





    \printbibliography

\end{document}